% ═══════════════════════════════════════════════════════════════
% PRÜFUNGSINFORMATION: Mündliche Prüfung Spanisch
% Klasse 10 Spätbeginner - Niveau A1
% ═══════════════════════════════════════════════════════════════

\documentclass[11pt, a4paper]{article}

\usepackage[utf8]{inputenc}
\usepackage[T1]{fontenc}
\usepackage[ngerman]{babel}

% --- SCHRIFTEN ---
\usepackage{helvet}
\renewcommand{\familydefault}{\sfdefault}
\usepackage{palatino}
\newcommand{\seriftext}[1]{{\fontfamily{ppl}\selectfont #1}}

\usepackage{microtype}
\usepackage[margin=2cm, top=2.2cm, bottom=1.8cm]{geometry}
\usepackage{enumitem}
\usepackage{tabularx}
\usepackage{booktabs}
\usepackage{array}
\usepackage{colortbl}
\usepackage{multicol}
\usepackage{tikz}
\usetikzlibrary{calc}

\usepackage{fancyhdr}
\usepackage{lastpage}

\usepackage{xcolor}
\usepackage{tcolorbox}
\tcbuselibrary{skins}

% ═══════════════════════════════════════════════════════════════
% FARBEN
% ═══════════════════════════════════════════════════════════════

\definecolor{spanischrot}{HTML}{C41E3A}
\definecolor{hellgrau}{HTML}{F5F5F5}
\definecolor{mittelgrau}{HTML}{666666}
\definecolor{dunkelgrau}{HTML}{333333}
\definecolor{rahmen}{HTML}{CCCCCC}

% ═══════════════════════════════════════════════════════════════
% BOXEN
% ═══════════════════════════════════════════════════════════════

\newtcolorbox{infobox}[1][]{
    colback=hellgrau, colframe=hellgrau, boxrule=0pt, arc=0pt,
    left=10pt, right=10pt, top=8pt, bottom=8pt,
    borderline west={3pt}{0pt}{spanischrot},
    before skip=8pt, after skip=8pt
}

\newtcolorbox{beispielbox}{
    colback=white, colframe=rahmen, boxrule=0.5pt, arc=0pt,
    left=10pt, right=10pt, top=8pt, bottom=8pt,
    before skip=6pt, after skip=6pt
}

\newtcolorbox{warnbox}{
    colback=white, colframe=spanischrot, boxrule=1pt, arc=0pt,
    left=10pt, right=10pt, top=8pt, bottom=8pt,
    before skip=8pt, after skip=8pt
}

% ═══════════════════════════════════════════════════════════════
% BEFEHLE
% ═══════════════════════════════════════════════════════════════

\newcommand{\es}[1]{\seriftext{\textbf{#1}}}
\newcommand{\teil}[2]{\noindent\textbf{\large Teil #1: #2}}

\setlength{\parindent}{0pt}
\setlength{\parskip}{6pt}

% ═══════════════════════════════════════════════════════════════
% KOPF- UND FUSSZEILE
% ═══════════════════════════════════════════════════════════════

\pagestyle{fancy}
\fancyhf{}
\renewcommand{\headrulewidth}{0.5pt}
\renewcommand{\footrulewidth}{0pt}
\lhead{\footnotesize\textsc{Spanisch} · Klasse 10}
\rhead{\footnotesize Mündliche Prüfung}
\cfoot{\footnotesize\thepage\,/\,\pageref{LastPage}}

% ═══════════════════════════════════════════════════════════════
% DOKUMENT
% ═══════════════════════════════════════════════════════════════

\begin{document}

% --- TITEL ---
\begin{center}
    {\LARGE\bfseries Prueba oral}\\[4pt]
    {\large Mündliche Prüfung Spanisch}\\[8pt]
    {\normalsize Klasse 10 · Spätbeginnerkurs · Niveau A1}
\end{center}

\vspace{8pt}

% ═══════════════════════════════════════════════════════════════
% ÜBERBLICK
% ═══════════════════════════════════════════════════════════════

\begin{infobox}
\textbf{Überblick}

\vspace{4pt}

\begin{tabular}{@{}ll@{}}
\textbf{Dauer:} & ca. 5--6 Minuten pro Person \\
\textbf{Hilfsmittel:} & keine \\
\textbf{Gesamtpunktzahl:} & 15 Punkte \\
\end{tabular}

\vspace{6pt}

Die Prüfung besteht aus \textbf{drei Teilen}:
\begin{enumerate}[leftmargin=2em, itemsep=2pt]
    \item \textbf{Producción oral} -- Monolog (1--2 Min) \dotfill 6 Punkte
    \item \textbf{Interacción} -- Fragen beantworten \dotfill 4 Punkte
    \item \textbf{Mediación} -- Sprachmittlung (30--60 Sek) \dotfill 3 Punkte
\end{enumerate}

\vspace{2pt}
Zusätzlich: \textbf{Aussprache \& Flüssigkeit} \dotfill 2 Punkte
\end{infobox}

\vspace{4pt}

% ═══════════════════════════════════════════════════════════════
% TEIL 1
% ═══════════════════════════════════════════════════════════════

\teil{1}{Producción oral -- Monolog} \hfill \textbf{6 Punkte}

\vspace{4pt}

Du sprichst \textbf{1--2 Minuten frei} auf Spanisch über dich selbst. Erzähle alles, was du über folgende Themen sagen kannst:

\begin{multicols}{2}
\begin{itemize}[leftmargin=1.5em, itemsep=2pt]
    \item Wie du heißt, woher du kommst, wo du wohnst
    \item Deine Familie (Eltern, Geschwister, Großeltern)
    \item Namen und Alter deiner Familienmitglieder
    \item Deine Persönlichkeit (Adjektive)
    \item Dein Stadtviertel/Wohnort (mit \es{hay}, \es{estar})
\end{itemize}
\end{multicols}

\begin{beispielbox}
\textbf{Beispielstart:}\\[4pt]
\es{¡Hola! Me llamo [Name] y tengo [X] años. Soy de Alemania y vivo en [Stadt]. En mi familia somos [X] personas: mi padre, mi madre y mi hermano/hermana...}\\[4pt]
\es{Mi padre se llama [Name] y tiene [X] años. Es muy simpático y...}\\[4pt]
\es{Vivo en un barrio tranquilo. Hay un parque, una biblioteca y...}
\end{beispielbox}

\vspace{6pt}

% ═══════════════════════════════════════════════════════════════
% TEIL 2
% ═══════════════════════════════════════════════════════════════

\teil{2}{Interacción -- Fragen beantworten} \hfill \textbf{4 Punkte}

\vspace{4pt}

Ich stelle dir \textbf{kurze Fragen} auf Spanisch. Du antwortest auf Spanisch. Die Fragen beziehen sich auf die gleichen Themen wie Teil 1 -- aber auch auf Dinge, die du vielleicht noch nicht erzählt hast.

\begin{beispielbox}
\textbf{Beispielfragen:}\\[4pt]
\begin{multicols}{2}
\begin{itemize}[leftmargin=1.5em, itemsep=3pt]
    \item \es{¿Cómo se llaman tus abuelos?}
    \item \es{¿Cuántos años tiene tu hermano/a?}
    \item \es{¿Vives en una ciudad o en un pueblo?}
    \item \es{¿Estudias o trabajas?}
    \item \es{¿Cómo es tu madre/padre?}
    \item \es{¿Qué hay en tu barrio?}
    \item \es{¿Dónde está la biblioteca?}
    \item \es{¿Tienes primos?}
\end{itemize}
\end{multicols}
\end{beispielbox}

\vspace{6pt}

% ═══════════════════════════════════════════════════════════════
% TEIL 3
% ═══════════════════════════════════════════════════════════════

\teil{3}{Mediación -- Sprachmittlung} \hfill \textbf{3 Punkte}

\vspace{4pt}

Ich sage dir einen \textbf{Satz auf Deutsch} und eine \textbf{Situation}. Du überträgst den Inhalt \textbf{sinngemäß ins Spanische}. Es geht \underline{nicht} um eine wörtliche Übersetzung!

\begin{beispielbox}
\textbf{Beispielaufgaben:}\\[4pt]
\begin{itemize}[leftmargin=1.5em, itemsep=4pt]
    \item \textit{Dein Freund ist neu in der Stadt. Erkläre ihm auf Spanisch, was es in deinem Stadtviertel gibt.}
    \item \textit{Erkläre auf Spanisch, wo die Bibliothek ist.}
    \item \textit{Beschreibe auf Spanisch, wie dein bester Freund / deine beste Freundin ist.}
    \item \textit{Erzähle auf Spanisch, was man in deinem Stadtviertel machen kann.}
\end{itemize}
\end{beispielbox}

\vspace{8pt}

% ═══════════════════════════════════════════════════════════════
% WICHTIGER HINWEIS
% ═══════════════════════════════════════════════════════════════

\begin{warnbox}
\textbf{Wichtig:} Die oben genannten Beispiele sind \underline{nicht erschöpfend}. Alle Inhalte, die wir im Unterricht behandelt haben, können in der Prüfung vorkommen. Bereite dich auf \textbf{alle Themen} vor!
\end{warnbox}

\newpage

% ═══════════════════════════════════════════════════════════════
% BEWERTUNGSKRITERIEN
% ═══════════════════════════════════════════════════════════════

\begin{center}
{\Large\bfseries Criterios de evaluación}\\[4pt]
{\normalsize Bewertungskriterien}
\end{center}

\vspace{8pt}

% --- TEIL 1: PRODUCCIÓN ORAL ---
\textbf{\large Teil 1: Producción oral (Monolog)} \hfill \textbf{max. 6 Punkte}

\vspace{4pt}

\begin{tabularx}{\textwidth}{|p{1.2cm}|X|p{1.5cm}|}
\hline
\rowcolor{hellgrau}
\textbf{Punkte} & \textbf{Beschreibung} & \textbf{Erreicht} \\
\hline
5--6 & Umfassende, zusammenhängende Darstellung. Alle Themenbereiche werden angesprochen. Gute Verwendung von Vokabular und Strukturen. Kaum Fehler. & \\
\hline
3--4 & Mehrere Themenbereiche werden angesprochen. Grundlegende Strukturen korrekt. Einige Fehler, die das Verständnis nicht beeinträchtigen. & \\
\hline
1--2 & Nur wenige Informationen. Häufige Fehler. Kommunikation teilweise eingeschränkt. & \\
\hline
0 & Keine zusammenhängende Äußerung möglich. & \\
\hline
\end{tabularx}

\vspace{12pt}

% --- TEIL 2: INTERACCIÓN ---
\textbf{\large Teil 2: Interacción (Fragen beantworten)} \hfill \textbf{max. 4 Punkte}

\vspace{4pt}

\begin{tabularx}{\textwidth}{|p{1.2cm}|X|p{1.5cm}|}
\hline
\rowcolor{hellgrau}
\textbf{Punkte} & \textbf{Beschreibung} & \textbf{Erreicht} \\
\hline
4 & Alle Fragen werden verstanden und angemessen beantwortet. Schnelle, natürliche Reaktionen. & \\
\hline
3 & Die meisten Fragen werden verstanden und beantwortet. Gelegentliches Nachfragen. & \\
\hline
2 & Etwa die Hälfte der Fragen wird verstanden/beantwortet. Häufigeres Nachfragen nötig. & \\
\hline
1 & Nur wenige Fragen werden verstanden. Häufige Missverständnisse. & \\
\hline
0 & Keine Kommunikation möglich. & \\
\hline
\end{tabularx}

\vspace{12pt}

% --- TEIL 3: MEDIACIÓN ---
\textbf{\large Teil 3: Mediación (Sprachmittlung)} \hfill \textbf{max. 3 Punkte}

\vspace{4pt}

\begin{tabularx}{\textwidth}{|p{1.2cm}|X|p{1.5cm}|}
\hline
\rowcolor{hellgrau}
\textbf{Punkte} & \textbf{Beschreibung} & \textbf{Erreicht} \\
\hline
3 & Inhalt wird vollständig und verständlich auf Spanisch übertragen. Angemessene Strukturen. & \\
\hline
2 & Inhalt wird größtenteils übertragen. Kleinere Auslassungen oder Ungenauigkeiten. & \\
\hline
1 & Nur Teile des Inhalts werden übertragen. Eingeschränkte Verständlichkeit. & \\
\hline
0 & Keine sinnvolle Übertragung möglich. & \\
\hline
\end{tabularx}

\vspace{12pt}

% --- AUSSPRACHE & FLÜSSIGKEIT ---
\textbf{\large Aussprache \& Flüssigkeit} \hfill \textbf{max. 2 Punkte}

\vspace{4pt}

\begin{tabularx}{\textwidth}{|p{1.2cm}|X|p{1.5cm}|}
\hline
\rowcolor{hellgrau}
\textbf{Punkte} & \textbf{Beschreibung} & \textbf{Erreicht} \\
\hline
2 & Klare, verständliche Aussprache. Natürlicher Redefluss mit angemessenen Pausen. & \\
\hline
1 & Aussprache meist verständlich. Einige Stockungen, aber Kommunikation nicht beeinträchtigt. & \\
\hline
0 & Aussprache stark beeinträchtigt. Sehr häufige Pausen/Abbrüche. & \\
\hline
\end{tabularx}

\vspace{16pt}

% --- GESAMTPUNKTZAHL ---
\begin{center}
\begin{tabular}{|p{8cm}|p{3cm}|}
\hline
\rowcolor{spanischrot}
\textcolor{white}{\textbf{GESAMTPUNKTZAHL}} & \textcolor{white}{\textbf{\_\_\_\_ / 15}} \\
\hline
\end{tabular}
\end{center}

\vspace{12pt}

% --- NOTENSCHLÜSSEL ---
\begin{center}
{\small
\begin{tabular}{|c|c|c|c|c|c|c|c|c|c|c|c|c|c|c|c|}
\hline
\rowcolor{hellgrau}
\textbf{NP} & 15 & 14 & 13 & 12 & 11 & 10 & 9 & 8 & 7 & 6 & 5 & 4 & 3 & 2 & 1 \\
\hline
\textbf{\%} & 100 & 96 & 91 & 86 & 80 & 73 & 67 & 60 & 53 & 47 & 40 & 33 & 27 & 20 & 13 \\
\hline
\textbf{P} & 15 & 14,4 & 13,6 & 12,9 & 12 & 11 & 10 & 9 & 8 & 7 & 6 & 5 & 4 & 3 & 2 \\
\hline
\end{tabular}}
\end{center}

\vspace{8pt}

\begin{center}
\textit{¡Buena suerte! -- Viel Erfolg!}
\end{center}

\end{document}
